\documentclass[a4paper, 14pt]{extarticle}

% Поля
%--------------------------------------
\usepackage{geometry}
\geometry{a4paper,tmargin=2cm,bmargin=2cm,lmargin=3cm,rmargin=1cm}
%--------------------------------------

\usepackage{minted}

%Russian-specific packages
%--------------------------------------
\usepackage[T2A]{fontenc}
\usepackage[utf8]{inputenc}
\usepackage[english, main=russian]{babel}

%--------------------------------------

\usepackage{textcomp}

% Красная строка
%--------------------------------------
\usepackage{indentfirst}               
%--------------------------------------             


%Graphics
%--------------------------------------
\usepackage{graphicx}
\graphicspath{ {./images/} }
\usepackage{float}%"Плавающие" картинки
%--------------------------------------

% Полуторный интервал
%--------------------------------------
\linespread{1.3}                    
%--------------------------------------

%Выравнивание и переносы
%--------------------------------------
% Избавляемся от переполнений
\sloppy
% Запрещаем разрыв страницы после первой строки абзаца
\clubpenalty=10000
% Запрещаем разрыв страницы после последней строки абзаца
\widowpenalty=10000
%--------------------------------------

%Списки
\usepackage{enumitem}

\usepackage{minted}
%Подписи
\usepackage{caption} 

%Гиперссылки
% \usepackage{hyperref}

\usepackage[colorlinks, urlcolor = blue, filecolor = blue, citecolor = blue, linkcolor = black]{hyperref}


\hypersetup {
unicode=true
}

%Рисунки
%--------------------------------------
\DeclareCaptionLabelSeparator*{emdash}{~--- }
\captionsetup[figure]{labelsep=emdash,font=onehalfspacing,position=bottom}
%--------------------------------------

\usepackage{tempora}

%Листинги
%--------------------------------------
%Листинги
%--------------------------------------
\usepackage{listings}
\usepackage{fancyvrb}
\usepackage{fvextra}
\usepackage{adjustbox}
\lstset{
  basicstyle=\ttfamily\scriptsize,
  breakatwhitespace=false,
  breaklines=false,
  captionpos=t,
  inputencoding=utf8,
  frame=single,
  keepspaces=true,
  columns=fullflexible,
  keywordstyle=\bfseries,
  numbers=left,
  numbersep=8pt,
  xleftmargin=18pt,
  xrightmargin=0pt,
  showspaces=false,
  showstringspaces=false,
  showtabs=false,
  stepnumber=1,
  tabsize=4,
  title=\lstname
}
%--------------------------------------

%%% Математические пакеты %%%
%--------------------------------------
\usepackage{amsthm,amsfonts,amsmath,amssymb,amscd}  % Математические дополнения от AMS
\usepackage{mathtools}                              % Добавляет окружение multlined
\usepackage[perpage]{footmisc}


% графики
\usepackage{booktabs}        % Пакет для улучшенных таблиц
\usepackage{pgfplots}        % Пакет для построения графиков
\pgfplotsset{compat=1.17}    % Установка совместимости
\usepackage{tikz}            % Пакет для создания диаграмм
\usetikzlibrary{arrows.meta, positioning, automata, decorations.pathreplacing, calc}

% Литература 

% переименовываем  список литературы в "список используемой литературы"
\addto\captionsrussian{\def\refname{Список литературы}}



\captionsetup[figure]{name=Рисунок, labelsep=space}


%--------------------------------------

%--------------------------------------
%			НАЧАЛО ДОКУМЕНТА
%--------------------------------------

\begin{document}

%--------------------------------------
%			ТИТУЛЬНЫЙ ЛИСТ
%--------------------------------------
\begin{titlepage}
\thispagestyle{empty}
\newpage


%Шапка титульного листа
%--------------------------------------
\vspace*{-60pt}
\hspace{-65pt}
\begin{minipage}{0.3\textwidth}
\hspace*{-20pt}\centering
\includegraphics[width=\textwidth]{emblem}
\end{minipage}
\begin{minipage}{0.67\textwidth}\small \textbf{
\vspace*{-0.7ex}
\hspace*{-6pt}\centerline{Министерство науки и высшего образования Российской Федерации}
\vspace*{-0.7ex}
\centerline{Федеральное государственное автономное образовательное учреждение }
\vspace*{-0.7ex}
\centerline{высшего образования}
\vspace*{-0.7ex}
\centerline{<<Московский государственный технический университет}
\vspace*{-0.7ex}
\centerline{имени Н.Э. Баумана}
\vspace*{-0.7ex}
\centerline{(национальный исследовательский университет)>>}
\vspace*{-0.7ex}
\centerline{(МГТУ им. Н.Э. Баумана)}}
\end{minipage}
%--------------------------------------

%Полосы
%--------------------------------------
\vspace{-25pt}
\hspace{-35pt}\rule{\textwidth}{2.3pt}

\vspace*{-20.3pt}
\hspace{-35pt}\rule{\textwidth}{0.4pt}
%--------------------------------------

\vspace{1.5ex}
\hspace{-35pt} \noindent \small ФАКУЛЬТЕТ\hspace{80pt} <<Информатика и системы управления>>

\vspace*{-16pt}
\hspace{47pt}\rule{0.83\textwidth}{0.4pt}

\vspace{0.5ex}
\hspace{-35pt} \noindent \small КАФЕДРА\hspace{50pt} <<Теоретическая информатика и компьютерные технологии>>

\vspace*{-16pt}
\hspace{30pt}\rule{0.866\textwidth}{0.4pt}

\vspace{11em}

\begin{center}
\Large {\bf Домашнее задание №1} \\ 
\large <<Имитационное моделирование систем массового обслуживания>> \\
\large {по курсу <<Моделирование>>} \\

\end{center}\normalsize

\vspace{8em}

\vbox{%
\hfill%
\vbox{%
\hbox{Студент: Аринин М.А.}%
\hbox{Группа: ИУ9-82Б}%
\hbox{Преподаватель: Домрачева А.Б.}%
}%
} 

\bigskip

\vfill


\begin{center}
\textsl{Москва 2026}
\end{center}
\end{titlepage}
%--------------------------------------
%		КОНЕЦ ТИТУЛЬНОГО ЛИСТА
%--------------------------------------


%\renewcommand{\ttdefault}{pcr}

\setlength{\tabcolsep}{3pt}

\newpage
\setcounter{page}{2}

\section{Постановка задачи}
Рассматривается задача имитационного моделирования процесса сдачи зачёта студентами первого курса вуза. Каждый студент имеет три попытки для сдачи зачёта по информатике с различными вероятностями успеха:
\begin{itemize}
  \item вероятность сдать с первой попытки: $p_1 = 0.6$;
  \item вероятность сдать со второй попытки: $p_2 = 0.25$;
  \item вероятность сдать с третьей попытки: $p_3 = 0.1$.
\end{itemize}

\textbf{Ограничения системы:}
\begin{itemize}
  \item между двумя попытками сдачи должно пройти не менее 1 дня;
  \item зачётная неделя длится 6 рабочих дней;
  \item в каждый из 6 рабочих дней может быть обслужено не более 20 попыток сдачи (остальные ждут в очереди).
\end{itemize}

\textbf{Цель работы:} провести вычислительный эксперимент с целью оценки количества студентов, успешно сдавших зачёт за 6 рабочих дней зачётной недели, для различных значений общего числа студентов ($N = 60, 80, 100$), используя систему имитационного моделирования GPSS.

%--------------------------------------------------------------------
\section{Построение математической модели}
Задача моделирования процесса сдачи зачёта может быть представлена как система массового обслуживания (СМО) с ограниченной пропускной способностью и вероятностными переходами между состояниями.

\subsection{Блок-схема процесса сдачи}
Каждый студент проходит через не более чем три попытки. После каждой неудачи студент возвращается в очередь на пересдачу через минимум 1 рабочий день, пока выполняются условия:
\begin{itemize}
  \item число попыток не превышает 3;
  \item попытка выполняется не позднее 6-го рабочего дня зачётной недели.
\end{itemize}

\begin{figure}[H]
\centering
\resizebox{\textwidth}{!}{%
\begin{tikzpicture}[
  >=Stealth,
  every node/.style={font=\small},
  box/.style={draw, rectangle, rounded corners, align=center},
  queueBox/.style={draw, rectangle, dashed, align=center, minimum width=3.9cm, minimum height=1.2cm},
  teacherBox/.style={draw, rectangle, align=center, minimum width=3.8cm, minimum height=1.3cm},
  successBox/.style={draw, rectangle, rounded corners, align=center, fill=green!18, minimum width=4.8cm, minimum height=0.9cm},
  fail3Box/.style={draw, rectangle, rounded corners, align=center, fill=red!12, text=red!60!black, minimum width=4.8cm, minimum height=0.9cm},
  failTimeBox/.style={draw, rectangle, rounded corners, align=center, fill=red!12, text=red!60!black, minimum width=4.8cm, minimum height=0.9cm},
  arr/.style={->, thick}
]
  % Блок: поступление студентов
  \node[box, minimum width=4.4cm, minimum height=1.1cm] (enter) at (0,0) {Входной поток\\студенты};

  % Очередь с "слотами"
  \node[queueBox] (queue) at (4.8,0) {};
  \node[font=\small] at ([yshift=0.24cm]queue.north) {Очередь};
  \node[font=\scriptsize] at ([yshift=-0.24cm]queue.south) {20 мест в день};
  % Рисуем условные слоты (просто чтобы визуально было как на скриншоте)
  \foreach \i in {0,...,4}{
    \draw[thick] (queue.south west) ++(0.55+\i*0.66,0.20) rectangle ++(0.40,0.60);
  }

  % Преподаватель
  \node[teacherBox] (teacher) at (10.05,0) {Преподаватель};

  % Ветки-результаты
  \node[successBox] (succ) at (14.85,1.4) {Успешная сдача};
  \node[fail3Box] (fail3) at (14.85,0.0) {Несдача за 3 попытки};
  \node[failTimeBox] (failTime) at (14.85,-1.45) {Не успели за 6 дней\\потратить 3 попытки};

  % Основные связи
  \draw[arr] (enter.east) -- (queue.west);
  \draw[arr] (queue.east) -- node[midway, above, font=\scriptsize] {сдача} (teacher.west);

  % Успех/неуспех
  \draw[arr] (teacher.east) -- (succ.west);
  \draw[arr] (teacher.east) -- (fail3.west);

  % Возврат на пересдачу (через преподавателя, как в схеме)
  \draw[arr] (teacher.south) to[bend left=35] node[midway, below=0pt, align=center] {Пересдача\\(число попыток $<3$)} (queue.south east);



  % Стрелка по ограничению времени от входного потока:
  % строго: вниз -> по низу вправо -> вверх в нижнюю грань блока "Не успели..."
  \coordinate (lateDown) at ($(enter.south)+(0,-2.05)$);
  \coordinate (lateTurn) at ($(failTime.south)+(0,-0.22)$);
  \draw[dashed, thick, ->]
    (enter.south) -- (lateDown)
    -- (lateDown -| lateTurn)
    -- (lateTurn)
    -- (failTime.south);

  % Прямая "скобка" под 90 градусов для выходных данных
  \coordinate (brTop) at ([xshift=0.45cm]succ.north east);
  \coordinate (brBot) at ([xshift=0.45cm]failTime.south east);
  \draw[thick] (brTop) -- ++(0.22,0);
  \draw[thick] ([xshift=0.22cm]brTop) -- ([xshift=0.22cm]brBot);
  \draw[thick] (brBot) -- ++(0.22,0);
  \node[align=left] at ([xshift=1.55cm]$(brTop)!0.5!(brBot)$) {Выходные\\данные};

\end{tikzpicture}
}
\caption{Схема процесса сдачи зачёта с очередью, пересдачами и ограничением времени}
\label{fig:flow}
\end{figure}

\subsection{Описание системы}
Система представляет собой СМО с следующими характеристиками:
\begin{itemize}
  \item \textbf{Входной поток:} студенты первого курса рассматриваются как транзакты модели. В базовом сценарии все студенты поступают в систему в день 1 и формируют общий поток заявок на сдачу.
  \item \textbf{Канал обслуживания:} обслуживание выполняет один преподаватель. За один рабочий день может быть обслужено не более 20 попыток сдачи, после достижения лимита оставшиеся транзакты переносятся на следующий день.
  \item \textbf{Время обслуживания:} длительность одной попытки в модели принимается пренебрежимо малой (результат попытки фиксируется мгновенно), однако между попытками одного и того же студента обязателен интервал не менее 1 рабочего дня.
  \item \textbf{Очередь:} очередь не ограничена по длине и формируется из студентов, ожидающих первую попытку, и студентов, направленных на пересдачу после неуспешной попытки.
  \item \textbf{Дисциплина обслуживания:} заявки из очереди обслуживаются по порядку поступления; повторная попытка добавляется в очередь как новая заявка после выдержки временного интервала.
  \item \textbf{Логика выхода из системы:} студент покидает систему либо при успешной сдаче, либо после трёх неуспешных попыток, либо при невозможности израсходовать все попытки в пределах 6 рабочих дней.
\end{itemize}


\subsection{Математическая модель вероятностей}
Для одного студента вероятность успешной сдачи зачёта вычисляется по формуле полной вероятности:
\[
P_{\text{успех}} = p_1 + (1-p_1) \cdot p_2 + (1-p_1) \cdot (1-p_2) \cdot p_3.
\]

Подставляя значения вероятностей:
\[
P_{\text{успех}} = 0.6 + 0.4 \cdot 0.25 + 0.4 \cdot 0.75 \cdot 0.1 = 0.6 + 0.1 + 0.03 = 0.73.
\]

Вероятность неуспешной сдачи:
\[
P_{\text{неуспех}} = (1-p_1) \cdot (1-p_2) \cdot (1-p_3) = 0.4 \cdot 0.75 \cdot 0.9 = 0.27.
\]

\subsection{Ограничения модели}
В реальной системе из-за ограничения в 6 дней и пропускной способности 20 попыток в день не все студенты могут использовать все три попытки. Это приводит к тому, что фактическая вероятность успеха может отличаться от теоретической $0.73$ из-за:
\begin{itemize}
  \item ограниченного времени (6 дней);
  \item ограниченной пропускной способности (20 попыток в день);
  \item зависимость результата от того, на какой день попадает очередная попытка (минимум 1 рабочий день между попытками).
\end{itemize}

\section{Решение вычислительных задач}
\label{sec:computational}

\subsection{Аналитическое решение (упрощенная модель)}
В упрощенной модели без учета временных и пропускных ограничений математическое ожидание количества успешно сдавших студентов для $N$ студентов:
\[
E[X] = N \cdot P_{\text{успех}} = N \cdot 0.73.
\]

Дисперсия количества успешно сдавших студентов (биномиальное распределение):
\[
\text{Var}[X] = N \cdot P_{\text{успех}} \cdot (1 - P_{\text{успех}}) = N \cdot 0.73 \cdot 0.27.
\]

Стандартное отклонение:
\[
\sigma = \sqrt{\text{Var}[X]} = \sqrt{N \cdot 0.73 \cdot 0.27}.
\]

\subsection{Учет ограничений системы}
В реальной системе с ограничениями необходимо учитывать:
\begin{enumerate}
  \item \textbf{Ограничение по времени:} студенты, начавшие сдавать в 3, 4, 5 и 6, не успевают сделать все три попытки.
  \item \textbf{Ограничение по пропускной способности:} если в день уже запланировано 20 попыток, новые студенты не могут попасть в этот день.
\end{enumerate}

Для точного аналитического решения с учетом всех ограничений требуется более сложный анализ, включающий:
\begin{itemize}
  \item распределение студентов по дням первой попытки;
  \item вероятности попадания на повторные попытки с учетом ограничений;
\end{itemize}

В данной работе для получения точной оценки используется имитационное моделирование в GPSS.

\section{Практическая реализация}
Модель реализована в системе имитационного моделирования GPSS World Student Version 5.2.
\subsection{Соответствие элементов модели операторам GPSS}
\begin{itemize}
  \item \textbf{Входной поток студентов:} оператор \texttt{GENERATE} в блоке \texttt{NEW\_ALL} формирует $N$ транзактов в момент времени $t=0$.
  \item \textbf{Параметры студента:} оператор \texttt{ASSIGN} использует \texttt{P1} для номера попытки и \texttt{P2} для дня первой попытки.
  \item \textbf{Ограничение 20 попыток в день:} проверка \texttt{TEST} по \texttt{X\$DAYUSED} и инкремент счётчика через \texttt{SAVEVALUE DAYUSED+,1}.
  \item \textbf{Смена дня и календарь:} блок \texttt{DAY\_START} обновляет \texttt{CURDAY} и сбрасывает \texttt{DAYUSED}.
  \item \textbf{Очередь FIFO без приоритета:} все транзакты (новые и на пересдачу) возвращаются в один и тот же блок \texttt{WAIT\_SLOT}.
  \item \textbf{Вероятностный исход попытки:} оператор \texttt{TRANSFER} реализует вероятности успеха $0.6/0.25/0.1$ для 1-й, 2-й и 3-й попыток.
  \item \textbf{Интервал между попытками:} оператор \texttt{ADVANCE 1.01} обеспечивает минимум 1 рабочий день между пересдачами.
  \item \textbf{Итоговые счётчики:} \texttt{SAVEVALUE PASSED} фиксирует число успешно сдавших, завершение моделирования задаётся блоком \texttt{END\_SIM}.
\end{itemize}
\newenvironment{longlisting}{\captionsetup{type=listing}}{}
\begin{longlisting}
\caption{Модель сдачи зачёта студентами в GPSS}
\begin{minted}[frame=single,fontsize = \footnotesize, linenos, xleftmargin = 1.5em, breaklines=true]{text}
NSTUD      EQU 60
DAYCAP     EQU 20
MAXDAY     EQU 6

SIMULATE

INIT        GENERATE 0,0,0,1
            SAVEVALUE PASSED,0
            SAVEVALUE DAYUSED,0
            SAVEVALUE CURDAY,1
            TERMINATE

NEW_ALL     GENERATE 0,0,0,NSTUD
            ASSIGN 1,1
            ASSIGN 2,0
            TRANSFER ,WAIT_SLOT

WAIT_SLOT   TEST LE X$CURDAY,MAXDAY,TIME_OVER
            TEST L X$DAYUSED,DAYCAP,WAIT_NEXT
            TRANSFER ,ENTER_EXAM

WAIT_NEXT   ADVANCE 0.01
            TRANSFER ,WAIT_SLOT

ENTER_EXAM  SAVEVALUE DAYUSED+,1
            TEST E P2,0,NOT_FIRST

FIRST_ATT   ASSIGN 2,X$CURDAY
            TRANSFER ,ATTEMPT_1

NOT_FIRST   TEST E P1,2,CHECK_ATT_3
            TRANSFER ,ATTEMPT_2

CHECK_ATT_3 TEST E P1,3,BAD_ATTEMPT
            TRANSFER ,ATTEMPT_3

BAD_ATTEMPT TERMINATE

ATTEMPT_1   TRANSFER 0.6,FAILED_1,PASSED_OK
ATTEMPT_2   TRANSFER 0.25,FAILED_2,PASSED_OK
ATTEMPT_3   TRANSFER 0.1,FINAL_FAIL,PASSED_OK

PASSED_OK   SAVEVALUE PASSED+,1
            TERMINATE

FAILED_1    TEST L X$CURDAY,MAXDAY,NO_RETRY_2
            ASSIGN 1,2
            ADVANCE 1.01
            TRANSFER ,WAIT_SLOT

NO_RETRY_2  TERMINATE

FAILED_2    TEST L X$CURDAY,MAXDAY,NO_RETRY_3
            ASSIGN 1,3
            ADVANCE 1.01
            TRANSFER ,WAIT_SLOT

NO_RETRY_3  TERMINATE

FINAL_FAIL  TERMINATE
TIME_OVER   TERMINATE

DAY_START   GENERATE 1,0,1,5
            SAVEVALUE DAYUSED,0
            SAVEVALUE CURDAY+,1
            TERMINATE

END_SIM     GENERATE 0,0,MAXDAY,1
            TERMINATE 1
\end{minted}
\end{longlisting}

\section{Результаты}
Для иллюстрации приведен полный фрагмент отчета GPSS World для эксперимента с $N=60$:

\begin{longlisting}
\caption{Отчет GPSS World Simulation Report для эксперимента $N=60$}
\begin{minted}[frame=single,fontsize = \footnotesize, linenos, xleftmargin = 1.5em, breaklines=true]{text}
START TIME           END TIME  BLOCKS  FACILITIES  STORAGES
     0.000              6.000    46        0          0

NAME                       VALUE  
ATTEMPT_1                  24.000
ATTEMPT_2                  25.000
ATTEMPT_3                  26.000
BAD_ATTEMPT                23.000
CHECK_ATT_3                21.000
CURDAY                  10005.000
DAYCAP                     20.000
DAYUSED                 10004.000
DAY_START                  41.000
END_SIM                    45.000
ENTER_EXAM                 15.000
FAILED_1                   29.000
FAILED_2                   34.000
FINAL_FAIL                 39.000
FIRST_ATT                  17.000
INIT                        1.000
MAXDAY                      6.000
NEW_ALL                     6.000
NOT_FIRST                  19.000
NO_RETRY_2                 33.000
NO_RETRY_3                 38.000
NSTUD                      60.000
PASSED                  10003.000
PASSED_OK                  27.000
TIME_OVER                  40.000
WAIT_NEXT                  13.000
WAIT_SLOT                  10.000

LABEL              LOC  BLOCK TYPE     ENTRY COUNT CURRENT COUNT RETRY
INIT                1    GENERATE             1             0       0
                    2    SAVEVALUE            1             0       0
                    3    SAVEVALUE            1             0       0
                    4    SAVEVALUE            1             0       0
                    5    TERMINATE            1             0       0
NEW_ALL             6    GENERATE            60             0       0
                    7    ASSIGN              60             0       0
                    8    ASSIGN              60             0       0
                    9    TRANSFER            60             0       0
WAIT_SLOT          10    TEST              8098             0       0
                   11    TEST              8098             0       0
                   12    TRANSFER            98             0       0
WAIT_NEXT          13    ADVANCE           8000             0       0
                   14    TRANSFER          8000             0       0
ENTER_EXAM         15    SAVEVALUE           98             0       0
                   16    TEST                98             0       0
FIRST_ATT          17    ASSIGN              60             0       0
                   18    TRANSFER            60             0       0
NOT_FIRST          19    TEST                38             0       0
                   20    TRANSFER            20             0       0
CHECK_ATT_3        21    TEST                18             0       0
                   22    TRANSFER            18             0       0
BAD_ATTEMPT        23    TERMINATE            0             0       0
ATTEMPT_1          24    TRANSFER            60             0       0
ATTEMPT_2          25    TRANSFER            20             0       0
ATTEMPT_3          26    TRANSFER            18             0       0
PASSED_OK          27    SAVEVALUE           43             0       0
                   28    TERMINATE           43             0       0
FAILED_1           29    TEST                20             0       0
                   30    ASSIGN              20             0       0
                   31    ADVANCE             20             0       0
                   32    TRANSFER            20             0       0
NO_RETRY_2         33    TERMINATE            0             0       0
FAILED_2           34    TEST                18             0       0
                   35    ASSIGN              18             0       0
                   36    ADVANCE             18             0       0
                   37    TRANSFER            18             0       0
NO_RETRY_3         38    TERMINATE            0             0       0
FINAL_FAIL         39    TERMINATE           17             0       0
TIME_OVER          40    TERMINATE            0             0       0
DAY_START          41    GENERATE             5             0       0
                   42    SAVEVALUE            5             0       0
                   43    SAVEVALUE            5             0       0
                   44    TERMINATE            5             0       0
END_SIM            45    GENERATE             1             0       0
                   46    TERMINATE            1             0       0

SAVEVALUE               RETRY       VALUE
 PASSED                   0         43.000                            
 DAYUSED                  0              0                            
 CURDAY                   0          6.000                            
\end{minted}
\end{longlisting}

\section{Анализ модели и проверка качества}
\subsection{Анализ результатов}
По результатам трех запусков модели (для $N=60, 80, 100$) получены следующие агрегированные показатели.

\begin{table}[H]
\centering
\begin{tabular}{|c|c|c|c|c|}
\hline
\textbf{$N$} & \textbf{ATTEMPT\_1} & \textbf{ATTEMPT\_2} & \textbf{ATTEMPT\_3} & \textbf{PASSED} \\
\hline
60 & 60 & 20 & 18 & 43 \\
\hline
80 & 80 & 30 & 10 & 52 \\
\hline
100 & 100 & 11 & 9 & 60 \\
\hline
\end{tabular}
\caption{Сводные результаты экспериментов}
\label{tab:results}
\end{table}

\textbf{Наблюдения по сериям экспериментов:}
\begin{itemize}
  \item Для $N=60$ система не насыщается по лимиту 120 попыток: выполнено 98 попыток, из них 43 успешные сдачи.
  \item Для $N=80$ суммарно выполнено 120 попыток, дневной лимит исчерпывается, часть студентов не успевает дойти до повторных попыток.
  \item Для $N=100$ также достигается верхний предел 120 попыток, при этом число вторых и третьих попыток снижается из-за роста очереди.
\end{itemize}

\textbf{Доли успешно сдавших по экспериментам:}
\begin{itemize}
  \item $N=60$: $\frac{43}{60}\approx 71.7\%$;
  \item $N=80$: $\frac{52}{80}=65.0\%$;
  \item $N=100$: $\frac{60}{100}=60.0\%$.
\end{itemize}

\subsection{Проверка качества модели}
\textbf{Проверка соответствия постановке задачи:}
\begin{itemize}
  \item Вероятности исходов по попыткам соответствуют условию: $0.6$, $0.25$, $0.1$.
  \item В модели реализовано ограничение не более 20 попыток в день (по счётчику \texttt{DAYUSED}).
  \item Между попытками одного студента соблюдается интервал не менее 1 дня (\texttt{ADVANCE 1.01}).
  \item Поведение очереди соответствует FIFO без приоритета пересдач.
\end{itemize}

\textbf{Проверка на граничных режимах:}
\begin{itemize}
  \item Для $N=60$ лимит 120 попыток не исчерпывается (выполнено 98 попыток).
  \item Для $N=80$ и $N=100$ достигается верхняя граница 120 попыток, что подтверждает корректную работу ограничения по дневной пропускной способности.
  \item Наличие значений в \texttt{NO\_RETRY\_2}, \texttt{NO\_RETRY\_3} и очереди на конец недели согласуется с ограничением в 6 рабочих дней.
\end{itemize}


\section{Выводы}
\begin{itemize}
  \item Выполнена формализация задачи имитационного моделирования процесса сдачи зачёта студентами как системы массового обслуживания с ограниченной пропускной способностью и вероятностными переходами между состояниями.
  
  \item Построена и реализована модель GPSS с общей FIFO-очередью без приоритета пересдач, с ограничением не более 20 попыток в день и минимальным интервалом 1 рабочий день между попытками одного студента.
  
  \item Проведены вычислительные эксперименты для $N=60$, $N=80$ и $N=100$ студентов; получены значения успешно сдавших 43, 52 и 60 соответственно.
  
  \item Сопоставление с упрощенной аналитической оценкой $E[X]=0.73N$ показало, что ограничение по числу попыток в день и конечная длительность зачётной недели существенно влияют на итоговый результат, особенно при росте входного потока.
  
\end{itemize}

\end{document}

